\section{Limitations}\label{sec:limitations}

The evidence is limited first by its data source. All models were developed and evaluated on governed partitions of one public SGCC dataset. Customer-level separation of development, confirmation, and final roles reduces overlap within this source, but it does not test transfer to another utility, geography, meter population, acquisition period, tariff regime, or inspection process. Published studies using the SGCC name are not automatically protocol-matched: preprocessing, balancing, prevalence, split membership, selection procedures, evaluation units, and metric definitions vary. The present results therefore support comparisons among the frozen AMP models and controls, not an unqualified ranking against published methods or a claim of external generalisation.

The annotation granularity imposes a second, more fundamental boundary. SGCC supplies one binary \texttt{FLAG} per customer and no verified theft-event timestamps, durations, or theft-containing windows. The 30-day windows used by AMP-LDP are measurement segments whose individual labels are unknown. Consequently, the model establishes customer-level discrimination only. Neither its latent coordinates nor a customer's high predicted probability identifies which window contains theft, and the study does not validate theft-event localisation or individual-window theft probabilities. Any event-level interpretation would require independently verified temporal annotations rather than propagation of the customer flag.

\textbf{Final evaluation cohort with disclosed historical baseline access.} This cohort had been accessed in Stage~13 for an obsolete supervised window-level baseline. The AMP-LDP hypothesis, architecture, distribution pooling, seed ensemble, controls, and confirmation decision were developed and frozen later without Study~2 row-level access to it. Stage~27A then fixed the no-refit policy, model hashes, comparison rules, uncertainty procedure, and outcome criteria before the terminal evaluation. This governance protects the prospective Study~2 decision chain, but it does not erase the earlier project-level access. The cohort must therefore be treated as a governed final evaluation with a disclosed limitation, not as a pristine, untouched, unseen, or never-accessed test set.

The cross-cohort summary also contains a reporting asymmetry. On \textsc{dev\_val}, the recorded AMP-LDP values are the arithmetic means of five individually computed seed metrics. The confirmation and final values are instead metrics of the arithmetic mean of the five customer-probability vectors. No separate frozen aggregate artifact stores a development ensemble point metric, so it was not regenerated. Directional consistency across the three roles is informative, but their AMP-LDP point estimates are not a strictly like-for-like ensemble trajectory, and changes in effect magnitude should not be attributed solely to cohort differences.

The method is conditional on a specific representation design: complete non-overlapping 30-day windows, the handcrafted \textsc{summary\_17} features, a $17\rightarrow32\rightarrow16$ encoder, six selected distribution statistics, HGB classification, and five seeds. The completed experiments do not establish that 30 days is the optimal window length, that the summaries are sufficient, or that the chosen quantiles and permutation-invariant aggregation are universally preferable. FULL\_48 also includes engineered chronology and persistence information absent from AMP-LDP and DIRECT; it is a strong benchmark rather than a representation-identical control. The more tightly matched DIRECT comparison holds aggregation, classifier configuration, and training population constant, but it still supports a representational association within this design rather than causal proof about a unique neural component.

The primary claim applies to the preregistered five-seed ensemble. It was numerically stronger than every individual seed on the final cohort, while four of five individual seeds met both frozen FULL\_48 parity margins. This passed the preregistered no-collapse rule, but it does not imply that an arbitrary single run will reproduce the ensemble result. The ensemble also requires five encoder--classifier pipelines rather than one. That structure may matter operationally, but final inference latency, throughput, memory use, energy consumption, model maintenance burden, and hardware sensitivity were not benchmarked, so no deployment-cost conclusion is warranted.

Finally, the study evaluates ranking through customer-level ROC-AUC and average precision. It does not establish probability calibration, a utility-specific intervention threshold, inspection yield under capacity constraints, economic benefit, robustness to future distribution shift, or causal explanations of theft behaviour. No field deployment or prospective operational trial was conducted. These are demonstrated absences in the evidence; possible effects of different utilities, changing consumption patterns, or operational constraints remain plausible concerns rather than observed failures of AMP-LDP.
